% GAIN 的两层图：左为异构提及图（同实体边、同句边、文档边），右为实体级图与两跳路径。
\begin{tikzpicture}[
  x=1mm,y=1mm,figure picture,
  mention/.style={figure node,minimum height=9mm,minimum width=13mm,
    draw=FigureInput,fill=FigureInputFill,font=\figurefont\scriptsize},
  docnode/.style={figure node,minimum height=7mm,minimum width=14mm,
    draw=FigureModel,font=\figurefont\scriptsize},
  entity/.style={figure node,minimum height=9mm,minimum width=15mm,
    draw=FigureInput,fill=FigureInputFill,font=\figurefont\scriptsize},
  intra/.style={draw=FigureStroke,line width=1.1pt},
  inter/.style={draw=FigureStroke,line width=.8pt,dash pattern=on 3pt off 1.5pt},
  docedge/.style={draw=FigureMuted,line width=.6pt,dash pattern=on .8pt off 1.2pt},
  path/.style={-{Stealth[length=1.8mm,width=1.2mm]},draw=FigureOutput,line width=1.1pt},
  query/.style={draw=FigureMuted,line width=.8pt,dash pattern=on 2pt off 2pt},
  note/.style={font=\figurefont\scriptsize,text=FigureMuted,align=center}
]
  \figurePanel{0}{0}{80}{48}{异构提及图}{FigurePanel}
  \figurePanel{88}{0}{60}{48}{实体级图与路径推理}{FigurePanel}

  % 提及图
  \node[docnode] (d) at (40,32) {文档节点};
  \node[mention] (m1) at (10,15) {$m_1$\\海棠大学};
  \node[mention] (m2) at (30,15) {$m_2$\\机器学习\\中心};
  \node[mention] (m3) at (50,15) {$m_3$\\该中心};
  \node[mention] (m4) at (70,15) {$m_4$\\读书会};
  \foreach \m in {m1,m2,m3,m4}{\draw[docedge] (d) -- (\m);}
  \draw[inter] (m1) -- (m2);
  \draw[intra] (m2) -- (m3);
  \draw[inter] (m3) -- (m4);
  \node[note] at (20,4.5) {句$S_1$};
  \node[note] at (60,4.5) {句$S_2$};

  % 实体图
  \node[entity] (e1) at (100,14) {$e_1$\\海棠大学};
  \node[entity] (e2) at (118,30) {$e_2$\\机器学习中心};
  \node[entity] (e3) at (136,14) {$e_3$\\读书会};
  \draw[path] (e1) -- (e2);
  \draw[path] (e2) -- (e3);
  \draw[query] (e1) -- node[figure edge label,font=\figurefont\scriptsize] {?} (e3);
  \node[note] at (118,4.5) {目标实体对$(e_1,e_3)$经$e_2$的两跳路径};
\end{tikzpicture}
